\section{System Architecture}
\label{sec:system}

Fig.~\ref{fig:arch} shows the Graphify pipeline. A software corpus is ingested
through five stages: \emph{detect}, \emph{extract} (two passes), \emph{merge},
\emph{cluster}, and \emph{query}. We describe each stage in turn.

\begin{figure*}[t]
  \centering
  % Architecture diagram for Graphify pipeline
% Full-width (figure*), two-column IEEE layout
\begin{tikzpicture}[
  font=\small,
  node distance=0.55cm and 0.7cm,
  % Box styles
  corpus/.style={draw, rounded corners=3pt, fill=blue!8, minimum width=1.6cm,
                 minimum height=0.55cm, align=center, font=\scriptsize},
  stage/.style={draw, rounded corners=3pt, fill=orange!12, minimum width=1.7cm,
                minimum height=0.65cm, align=center, font=\scriptsize\bfseries},
  pass/.style={draw, rounded corners=3pt, fill=green!10, minimum width=2.0cm,
               minimum height=0.55cm, align=center, font=\scriptsize},
  output/.style={draw, rounded corners=3pt, fill=purple!10, minimum width=1.7cm,
                 minimum height=0.55cm, align=center, font=\scriptsize},
  platform/.style={draw, rounded corners=3pt, fill=gray!12, minimum width=1.5cm,
                   minimum height=0.5cm, align=center, font=\scriptsize},
  arrow/.style={-{Stealth[length=4pt]}, thick},
  tag/.style={font=\scriptsize\itshape, text=gray},
]

% ---- CORPUS (left column) ----
\node[corpus] (code)   at (0, 0)    {Code\\(33 langs)};
\node[corpus] (docs)   [below=0.2cm of code]   {Docs / MD};
\node[corpus] (pdfs)   [below=0.2cm of docs]   {Papers / PDF};
\node[corpus] (images) [below=0.2cm of pdfs]   {Images};
\node[corpus] (video)  [below=0.2cm of images] {Video / Audio};
\node[corpus] (office) [below=0.2cm of video]  {Office (.docx)};

% label
\node[tag, above=0.15cm of code] {Corpus};

% ---- DETECT ----
\node[stage] (detect) [right=1.1cm of docs, yshift=-0.5cm] {Detect \&\\Classify};

% corpus -> detect arrows
\foreach \src in {code, docs, pdfs, images, video, office}
  \draw[arrow, gray!60] (\src.east) -- (detect.west |- \src.east);

% ---- EXTRACT (two passes) ----
\node[pass] (ast) [right=1.0cm of detect, yshift=0.6cm]
      {Pass 1:\\AST (tree-sitter)\\{\color{green!60!black}\scriptsize EXTRACTED}};
\node[pass] (llm) [right=1.0cm of detect, yshift=-0.6cm]
      {Pass 2:\\LLM Subagent\\{\color{orange!70!black}\scriptsize INFERRED}};

\draw[arrow] (detect.east) -- ++(0.35,0) |- (ast.west);
\draw[arrow] (detect.east) -- ++(0.35,0) |- (llm.west);

% label
\node[tag, above=0.15cm of ast] {Extract};

% ---- MERGE ----
\node[stage] (merge) [right=1.0cm of ast, yshift=-0.6cm] {Merge \&\\Dedup};
\draw[arrow] (ast.east) -- ++(0.35,0) |- (merge.west);
\draw[arrow] (llm.east) -- ++(0.35,0) |- (merge.west);

% ---- CLUSTER ----
\node[stage] (cluster) [right=1.0cm of merge] {Leiden\\Cluster};
\draw[arrow] (merge) -- (cluster);
\node[tag, above=0.1cm of cluster] {Cluster};

% ---- GRAPH.JSON ----
\node[output] (graph) [right=1.0cm of cluster] {\texttt{graph.json}\\{\scriptsize nodes, edges,}\\{\scriptsize communities}};
\draw[arrow] (cluster) -- (graph);

% ---- QUERY ENGINE ----
\node[stage] (query) [below=0.8cm of graph] {Hub-Aware\\BFS/DFS};
\draw[arrow] (graph.south) -- (query.north);
\node[tag, right=0.05cm of query] {\scriptsize $\tau = p_{99}(\deg)$};

% user query input
\node[tag] (qin) [left=0.6cm of query] {\textit{query $q$}};
\draw[arrow, dashed] (qin) -- (query.west);

% subgraph output
\node[output] (subgraph) [below=0.5cm of query] {Scoped\\Subgraph};
\draw[arrow] (query) -- (subgraph);

% ---- PLATFORMS ----
\node[platform] (p1) [right=1.0cm of graph, yshift=0.4cm]  {Claude Code};
\node[platform] (p2) [right=1.0cm of graph, yshift=-0.2cm] {Cursor};
\node[platform] (p3) [right=1.0cm of graph, yshift=-0.8cm] {Gemini CLI};
\node[platform] (p4) [right=1.0cm of graph, yshift=-1.4cm] {+18 more};

\foreach \p in {p1, p2, p3, p4}
  \draw[arrow, gray!70] (graph.east) -- (\p.west);

\node[tag, above=0.1cm of p1] {21 Platforms};

% ---- HOOKS annotation ----
\node[tag] (hooks) [below=0.3cm of p4] {\scriptsize PreToolUse hooks};
\node[tag] [below=0.05cm of hooks] {\scriptsize + Progressive skills};

\end{tikzpicture}

  \caption{Graphify pipeline. A heterogeneous corpus (code, docs, PDFs, images,
           video) flows through multi-modal detection, hybrid two-pass extraction,
           graph merge and deduplication, topology-only community detection, and
           hub-aware traversal. Every edge carries a confidence tag.
           The graph is exposed to 21 AI assistant platforms via always-on hooks
           and progressive-disclosure skill files.}
  \label{fig:arch}
\end{figure*}

\subsection{Multi-Modal Detection and Classification}

Graphify classifies each file into one of six types: \textsc{Code},
\textsc{Document}, \textsc{Paper}, \textsc{Image}, \textsc{Video}, or
\textsc{Office}. Classification is based on file extension against curated
extension sets: 80+ patterns for code (covering 33 languages), plus PDF,
image, video, and office formats. An incremental manifest records file
modification timestamps; only files whose \texttt{mtime} changed since the
last extraction are re-processed, enabling sub-second updates on typical commits.

Security measures are applied before parsing untrusted corpus files. Office
documents (\texttt{.docx}, \texttt{.xlsx}) are screened via a bounded
streaming decompression check: if the total decompressed size of ZIP members
exceeds $D_{\max} = 512\,\text{MiB}$ or the compression ratio exceeds $\rho =
200$, the file is skipped. PDF files are subject to a raw-size cap of
$50\,\text{MiB}$. URL ingestion is protected by a SSRF guard that resolves
hostnames and rejects private, link-local, and cloud-metadata IP ranges.

\subsection{Hybrid Two-Pass Extraction}

\subsubsection{Pass 1: Deterministic AST Extraction}

For each \textsc{Code} file, Graphify invokes the language-specific
tree-sitter~\cite{treesitter} grammar to parse the source into a concrete
syntax tree. A language-aware walker then emits nodes and edges:

\begin{itemize}
  \item \textbf{Nodes}: function definitions, class declarations, variable
        bindings, and file-level module nodes. Each node carries a
        \texttt{source\_location} stamp (e.g., \texttt{L42}).
  \item \textbf{Edges}: \texttt{calls}, \texttt{imports}, \texttt{contains},
        \texttt{inherits}, \texttt{implements}, \texttt{embeds}.
\end{itemize}

All AST-derived edges receive confidence \textsc{Extracted} ($c = 1.0$),
reflecting that the relation is syntactically explicit in the source. A
language-specific built-in filter suppresses universal identifiers
(\texttt{print}, \texttt{len}, \texttt{String}, \texttt{Promise}) that would
otherwise become pathological god nodes. Import resolution is language-aware:
Python relative imports are resolved against the module hierarchy; TypeScript
\texttt{tsconfig} path aliases are respected; Go module paths are normalized.

\subsubsection{Pass 2: LLM Semantic Extraction}

For non-code files (\textsc{Document}, \textsc{Paper}, \textsc{Image},
\textsc{Video}, \textsc{Office}), Graphify dispatches an LLM subagent.
Video and audio files are first transcribed to text via faster-whisper~\cite{whisper};
PDF files are parsed via \texttt{pypdf}; images are passed directly to a
vision-capable model. The subagent receives a structured extraction prompt
specifying the node-ID format, relation taxonomy, confidence rubric, and output
JSON schema. It emits nodes and edges with one of three confidence levels:

\begin{align}
  c(e) &= \begin{cases}
    1.0 & \text{\textsc{Extracted}: explicit in source} \\
    [0.75,\,1.0) & \text{\textsc{Inferred}: LLM-reasoned} \\
    [0,\,0.75) & \text{\textsc{Ambiguous}: low-confidence, flagged}
  \end{cases}
  \label{eq:confidence}
\end{align}

The semantic pass is backend-agnostic: Claude Code subagents (default), OpenAI,
Gemini, Ollama, and any OpenAI-compatible endpoint are supported.

\subsubsection{Graph Merge and Deduplication}

Nodes from both passes are merged idempotently into a NetworkX directed graph.
On ID collision, the semantic node overwrites the AST node, preserving the richer
LLM-inferred metadata. Three deduplication layers are applied: (1) a
within-file \texttt{seen\_ids} set collapses duplicate definitions; (2) the
graph's \texttt{add\_node} semantics overwrite stale entries from previous
incremental runs; (3) an optional LLM tiebreaker resolves near-duplicate nodes
identified via MinHash LSH~\cite{minhash} blocking and
Jaro-Winkler similarity~\cite{jaro}.

\subsubsection{Graph Schema}

The unified schema is:
{\small
\begin{align*}
  \text{Node} &= \langle \mathtt{id},\; \mathtt{label},\; \mathtt{src\_file},\;
  \mathtt{loc},\; \mathtt{type},\; \mathtt{community} \rangle \\
  \text{Edge} &= \langle \mathtt{src},\; \mathtt{tgt},\; \mathtt{relation},\;
  c,\; \mathtt{weight} \rangle
\end{align*}
}

\noindent where \texttt{relation} $\in \{$\texttt{calls, imports, contains,
inherits, implements, uses, references,}\\
\texttt{cites, transcribes, embeds}$\}$.
This schema is modality-agnostic: a Python function and a paragraph in a PDF
design document are both first-class nodes; a \texttt{references} edge connects
them. The graph is serialized as NetworkX node-link JSON, which is git-merge-safe
and diffable line-by-line.

\subsection{Embedding-Free Community Detection}

Graphify partitions the graph into communities using the Leiden
algorithm~\cite{leiden}, which optimizes modularity:
\begin{equation}
  Q = \frac{1}{2m} \sum_{ij} \left[ A_{ij} - \frac{k_i k_j}{2m} \right]
  \delta(c_i, c_j)
  \label{eq:modularity}
\end{equation}
where $m$ is the number of edges, $A_{ij}$ is the adjacency matrix, $k_i$ is
the degree of node $i$, and $\delta(c_i, c_j)$ is 1 iff $i$ and $j$ share a
community. Leiden~\cite{leiden} guarantees all communities are internally
connected and refines toward subset-optimal partitions, correcting Louvain's
known production of internally disconnected communities~\cite{louvain}.
Louvain via NetworkX is used as a fallback when graspologic is unavailable.
Note that community detection is applied to an undirected projection of the
graph, as modularity is defined for undirected graphs.

\medskip
\noindent\textbf{Hub exclusion.}
Nodes with degree $k_i \geq \tau$ where $\tau = \max(50,\; p_{99}(\{k_i\}))$
are excluded before running Leiden, then re-inserted into the community
holding the plurality of their neighbors (majority vote). This prevents
high-degree utility nodes (e.g., \texttt{logging}, \texttt{json.dumps})
from dominating community structure.

\medskip
\noindent\textbf{Determinism.}
Community IDs are assigned by total order $\langle -|C|, \text{sorted}(C)
\rangle$: larger communities receive smaller IDs; ties are broken
lexicographically by sorted node IDs. This ensures bit-for-bit identical
community assignments across runs on the same input.

\subsection{Hub-Aware Graph Traversal}
\label{subsec:traversal}

At query time, the user's natural language question $q$ is tokenized,
IDF-weighted, and matched against node labels. Each node $n$ is scored using
a three-tier match with a file-name bonus:
\begin{equation}
  s(n, q) = w_\text{IDF} \cdot \max\!\left(10^3 \cdot \mathbf{1}_\text{exact},\;
  10^2 \cdot \mathbf{1}_\text{prefix},\; \mathbf{1}_\text{substr}\right)
  + 0.5 \cdot \mathbf{1}_\text{file}
  \label{eq:score}
\end{equation}
where $w_\text{IDF}$ is the query-term IDF weight, and each indicator is 1
if the query term matches the node's label at the respective granularity
(exact, prefix, or substring); the max selects the highest-tier match.
The file bonus $0.5 \cdot \mathbf{1}_\text{file}$ fires when the query
matches the node's source filename. Top-$K$ seeds:
$S^* = \arg\!\max_{|S|=K} \sum_{n \in S} s(n, q)$.

Algorithm~\ref{alg:bfs} formalizes the hub-aware traversal:

\begin{algorithm}[h]
\caption{Hub-Aware BFS}
\label{alg:bfs}
\begin{algorithmic}[1]
\REQUIRE Graph $G=(V,E)$, seeds $S^* \subseteq V$, depth $d$, threshold $\tau$
\ENSURE Subgraph $H$
\STATE $\text{visited} \gets S^*$;\; $\text{frontier} \gets S^*$
\FOR{$i = 1$ \TO $d$}
  \STATE $\text{next} \gets \emptyset$
  \FOR{$n \in \text{frontier}$}
    \IF{$\deg(n) < \tau$}
      \STATE $\text{next} \gets \text{next} \cup (\mathcal{N}(n) \setminus \text{visited})$
    \ENDIF
  \ENDFOR
  \STATE $\text{visited} \gets \text{visited} \cup \text{next}$;\; $\text{frontier} \gets \text{next}$
\ENDFOR
\RETURN $G[\text{visited}]$
\end{algorithmic}
\end{algorithm}

Hub nodes (degree $\geq \tau$) are \emph{included} in the subgraph for context
but are \emph{not expanded}: their neighbors are not added to the frontier. This
prevents god-node explosion while preserving the hub's contextual relevance.

\subsection{AI Assistant Integration}

Graphify integrates with AI assistant platforms through two mechanisms:

\textbf{Always-on hooks.} On platforms supporting payload-bearing \texttt{PreToolUse}
hooks (Claude Code, Gemini CLI), two hooks are registered: one fires before
Bash search commands (\texttt{grep}, \texttt{rg}, \texttt{find}); the other
before native file-read tool calls (\texttt{Read}, \texttt{Glob}). Both hooks
check for the presence of \texttt{graphify-out/graph.json} and, when found,
inject an \texttt{additionalContext} message steering the agent toward
\texttt{graphify query} instead of raw file reading. Hooks never block: every
branch fails open, so legitimate reads always proceed.

\textbf{Progressive-disclosure skills.} Each platform receives a lean skill file
(615 lines) carrying the full default pipeline inline, plus an on-demand
\texttt{references/} sidecar (8 files, one per non-default workflow). An agent
reads a reference file only when its routing condition fires, reducing always-loaded
context by 47\% versus the prior monolithic 1,156-line skill file.


\subsection{Implementation and Availability}

Graphify is implemented in Python 3.10+ and published on PyPI as
\texttt{graphifyy} (MIT license). The system ships with tree-sitter grammars
for 33 language variants as compiled C extensions, requiring no external
toolchain. A single command --- \texttt{graphify extract .} --- runs the
full detection, extraction, clustering, and report generation pipeline.
The incremental update path (\texttt{graphify update .}) re-extracts only
modified files and is suitable for git post-commit hooks. The query interface
(\texttt{graphify query "<question>"}) is available both as a CLI command and
as an MCP (Model Context Protocol) server for programmatic integration.
Source code, benchmark data, and worked examples are available at:
\url{https://github.com/safishamsi/graphify}.
